\emph{函数}~$f\colon A \to B$ 将\emph{定义域}~$A$ 中的每个元素映射到\emph{陪域}~$B$ 中唯一的一个元素。
如果 $x \in A$，$f$ 将 $x$ 映射到的那个 $y$ 称为 $f$ 关于\emph{自变元}~$x$ 的\emph{值}~$f(x)$。
如果 $A$ 是一个序对的集合，我们就可以把函数~$f$ 看作接受两个自变元。
$f$ 的\emph{值域}~$\ran{f}$ 是~$B$ 的一个子集，由~$f$ 的所有值组成。

如果 $\ran{f} = B$，则称 $f$ 是\emph{满射的}。
值~$f(x)$ 的唯一性在于：$f$ 将 $x$ 只映射到一个 $f(x)$，绝不会多于一个。
如果 $f(x)$ 还满足这样一种唯一性，即任意两个不同的自变元都不会被映射到同一个值，则称 $f$ 是\emph{单射的}。
既是单射又是满射的函数称为\emph{双射的}。

双射函数有唯一的\emph{反函数}~$f^{-1}$。
函数也可以复合在一起：函数 $(g \circ f)$ 是 $f$ 与 $g$ 的\emph{复合}。
单射函数的复合仍是单射的，满射函数的复合仍是满射的，且 $(f^{-1} \circ f)$ 是恒等函数。

如果我们放宽"$f$ 必须对每个 $x \in A$ 都有值"这一要求，就得到了\emph{偏函数}的概念。
如果 $f\colon A \pto B$ 是偏函数，当 $f$ 对自变元~$x$ 有值时，我们说 $f(x)$ \emph{有定义}，记作 $f(x) \fdefined$；否则我们说 $f(x)$ \emph{无定义}，记作 $f(x) \fundefined$。
任何（偏）函数~$f$ 都有一个与之关联的\emph{图}~$R_f$，即当且仅当 $f(x) = y$ 时成立的关系。